\begin{abstract}
Ising machines solve combinatorial optimization problems by annealing a network of spins. On dense problems, where every spin couples to every other, the least expensive hardware step updates all spins at once: stochastic cellular automata (SCA) decide every spin in parallel and read couplings only for the spins that flip. Such synchronous updates, however, \textit{create an echo: the previous state of each spin returns to it through its neighbors}. On K2000, the Max-Cut benchmark on a complete graph of 2{,}000 nodes, two-thirds of all flips undo a flip that the same spin made one step earlier. This paper presents \emph{Anechoic}, an FPGA Ising machine that \textit{subtracts this echo from every decision}. Dynamical mean-field theory identifies the echo as an Onsager reaction and gives the strength of its one-step part: \textit{the number of spins whose flip probability lies strictly between 0 and 1, divided by twice the temperature}. The decision logic already flags these spins, so this Onsager correction costs one population count per step, and a temperature-scaled approximation needs no count. A hand-written RTL engine executes a dense 2{,}000-spin step in \cycStep{} cycles plus \cycFlip{} cycles per flip, so every step and flip that the correction saves is saved time. In experiments, the counted and the temperature-scaled correction both find better K2000 cuts at every step budget than each of five published binary-spin algorithms run at its paper's settings. Twelve engines at \fClk{} on an AMD Alveo V80 reach a cut of at least 33{,}000 on K2000 with 99\% probability in \textbf{\ttsPr{}, \spdBSB{} faster than the fastest published hardware}, at \ets{} per solution; with 2-bit couplings, the temperature-scaled form is also faster than uncorrected SCA on 49 of 51 sparse G-set graphs. On an NVIDIA GH200 GPU, the identical engine is \gpuSpd{} slower and needs \gpuEratio{} more energy per solution, though its 240 concurrent anneals give 2.3$\times$ higher throughput. As a GlobalFoundries 22\,nm ASIC routed at 1.25\,GHz (typical corner), twelve engines would reach the target in about \textbf{10\,\textmu s, 26$\times$ faster than the fastest published hardware}, with 0.20\,mJ of engine energy per solution. The code is available at \coderepo{}.
\end{abstract}
